\section{Experiments}
\label{sec:experiments}

The central question is not whether the causal PT beats a standard decoder, but whether
syntactic structure pays for itself. That question is answered by the gap to a Looped
Transformer, not by the gap to GPT (\Cref{sec:exp2}).

\NOTE{Every number in this section must be traceable to a row in the corresponding
\texttt{experiments/*/EXPERIMENT\_STATUS.md} run log --- date, commit, config, seed, metric,
wall-clock. Do not copy a number here that does not have such a row.}

\subsection{Shared setup}
\label{sec:setup}

Everything in \Cref{tab:setup} is held constant across models; any difference other than the
one under test would invalidate the comparison.

\begin{table}[t]
\centering
\begin{tabular}{ll}
\toprule
Item & Setting \\
\midrule
Data                   & \TODO{PTB or WikiText-2 --- state which, and why not WikiText-103} \\
Parameter budget       & \TODO{20--50M, matched; state the matching rule} \\
Tokenizer / vocabulary & identical across models \\
Context length         & identical across models \\
Training loop          & one shared implementation \\
Optimizer              & Adam, identical hyperparameters and schedule \\
Seeds                  & \TODO{$k$ seeds; report mean $\pm$ std} \\
Metric                 & perplexity on the held-out set \\
\bottomrule
\end{tabular}
\caption{Setup held constant across all models.}
\label{tab:setup}
\end{table}

\paragraph{On the corpus.}
WikiText-103 is deliberately excluded. \citet{wu2023probabilistic} report that the model
degrades beyond roughly $10^5$ sentences, suspected to be the absence of a feed-forward
structure; running there would test a known failure mode of the encoder rather than the causal
extension. \NOTE{say this before a reviewer asks why the corpus is small}

\paragraph{On matching budgets.}
We report parameter count \emph{and} wall-clock/FLOPs. The PT shares parameters across
iterations, so equal parameter count does not mean equal compute.

\paragraph{On the embedding split.}
We report embedding and non-embedding parameters separately, which matters more here than in a
usual comparison. With tied embeddings (\Cref{prop:tying}) the PT's parameters sit almost
entirely in $\Sfac \in \R^{|\Vocab| \times \nlab}$; the arc scores $\Tfac{c}$, the word unary
$\bfac$ and the root key $\rfac{c}$ are negligible beside it. At a matched total, the PT is
close to an embedding table plus a little, while a GPT baseline is not, and a single total
hides that completely. \TODO{state the matching rule chosen --- total parameters,
non-embedding parameters, or both with vocabulary fixed. Any is defensible; leaving it
unstated is what gets attacked first.}

\subsection{Experiment 0: implementation validation}
\label{sec:exp0}

Not a result, but the precondition for every result below, and reported because two of its
checks test the derivation rather than the code.

\begin{itemize}
  \item \textbf{Causality.} Changing token $t$ leaves the logits at every position $<t$
        bitwise unchanged (CPU, fixed seed).
  \item \textbf{Frozen prefix.} Prefix quantities are kept as attached leaves with
        \texttt{requires\_grad=True}; after \texttt{backward()} from the loss at step $t$ their
        gradient is zero (\Cref{sec:frozen-prefix}).
  \item \textbf{Tying.} Input and output embeddings are the same tensor object
        (\Cref{prop:tying}).
  \item \textbf{Free energy.} $\Free$ is non-increasing across mean-field iterations on a toy
        graph. \NOTE{this is the check that catches a typo in the update equations --- shapes,
        normalisation and causality are all blind to it}
  \item \textbf{Exact readout vs.\ brute force.} On one slot the graph is a tree, so
        \Cref{eq:exact-readout} must agree with explicit enumeration over a tiny vocabulary to
        numerical precision. This validates the factorisation itself.
  \item \textbf{Worked example.} Reproduces the numbers in \texttt{causal\_pt\_output\_note.pdf}
        \S5 ($\Vocab = \{\textit{the}, \textit{cat}, \textit{sat}, \textit{mat}\}$, $\nlab=2$,
        one channel, both readout modes).
\end{itemize}

\TODO{report pass/fail and the free-energy curve; a figure of $\Free$ vs.\ iteration is cheap
and is the most convincing single plot in this subsection}

\subsection{Experiment 1: causal PT vs.\ a standard decoder}
\label{sec:exp1}

\textbf{Question.} Does the causal PT train, and does its perplexity land in a reasonable
corridor? This is not a controlled comparison --- everything differs --- so it measures the
total cost of the construction, not the effect of any component.

\textbf{Arms.} Split into 1.1 (without the Appendix~B.3 global variables $\Gv{t}$) and 1.2
(with them). The globals are a \emph{measured} variable, not an assumption: \citet{wu2023probabilistic}
propose them as an in-graph substitute for the missing feed-forward structure but never test
them, so the delta between the arms is itself a result.

\TODO{table + training curves} \TODO{numbers}

\subsection{Experiment 2: causal PT vs.\ Looped Transformer --- core result}
\label{sec:exp2}

\textbf{Question.} Does syntactic structure contribute beyond weight sharing?

The PT differs from a standard decoder in two ways at once: it reuses one parameter set across
iterations, and it has a syntactic factor graph. A Looped Transformer --- one block applied
$\niter$ times --- has the first without the second (\Cref{tab:three-point}).

\begin{table}[t]
\centering
\begin{tabular}{lcc}
\toprule
Model & Weight sharing & Syntactic structure \\
\midrule
GPT                & no  & no  \\
Looped Transformer & yes & no  \\
Causal PT          & yes & yes \\
\bottomrule
\end{tabular}
\caption{Looped vs.\ GPT isolates weight sharing; PT vs.\ Looped isolates structure.
The second comparison is the claim of this paper.}
\label{tab:three-point}
\end{table}

\textbf{Control.} The loop depth $\niter$ of the Looped baseline matches the number of
mean-field iterations in the PT, so effective depth is comparable as well as parameter count.
We report a small sweep over $\niter$ rather than a single value, so the comparison does not
rest on a lucky choice.

\textbf{Falsification.} If the PT is worse than Looped at matched budget across the sweep, the
paper's central claim fails, and we report that. \NOTE{keep this sentence --- the run log is
required to keep failed and superseded rows, and the paper should be written by someone willing
to publish the negative result}

\TODO{table + numbers}

\subsection{Experiment 3: exact readout vs.\ mean-field readout}
\label{sec:exp3}

\textbf{Question.} What does the mean-field approximation cost? Following \Cref{sec:inference},
the exact readout is the reference and the mean-field one is the ablation.

\begin{description}
  \item[(a) Evaluation-time swap.] Train with the mean-field readout, substitute
        \Cref{eq:exact-readout} at evaluation on the same parameters. Isolates the
        approximation error of the readout. Nearly free.
  \item[(b) Train with each.] Two full runs; measures the end-to-end effect, including how the
        approximation shapes learning.
\end{description}

Neither outcome is a win or a loss: a small gap says mean-field is adequate and cheap, a large
one says the tree structure should be exploited. \TODO{also check whether the
mixture-of-softmaxes form relieves the rank-$\nlab$ bottleneck of \Cref{sec:bottleneck}}

\TODO{table + numbers}
